\chapter{متباينة هيلبرت ومدخل إلى طريقة القيم الكبيرة}
\personindex{هيلبرت، ديفيد}
\personindex{شور، إيساي}
\symbolindex{متباينة هيلبرت}

\section{نطاق الفصل وحالته}

\noindent\textbf{حالة الفصل:} \textenglish{AUTHORED-DRAFT / NON-CITABLE}.
يثبت هذا الفصل نتيجةً واحدة من الصفر: متباينة هيلبرت بثابت شور الحادّ
\(\pi\). أما مبرهنة القيم الكبيرة لكثيرات حدود ديريشليه — الأداة التي
تُستعمَل فعلًا في مسائل الكثافة الصفرية — فتُذكَر هنا \textbf{مُستشهَدًا
بها فقط}، ولا يُعاد بناء برهانها.

\begin{remark}[تحذير نطاق إلزامي: ما لا يفعله هذا الفصل]
\label{rem:ch58-scope-warning}
\resultid{ANT-PRIN-58-01}
\noindent\textenglish{METHODOLOGICAL-PRINCIPLE / INFERENCE-GUARDED}.
السلسلة المبرهَنة في هذا الفصل \textbf{لا تكفي بذاتها} لاشتقاق أي حدّ
كثافة صفرية. ما يُثبَت هنا هو الحالة الكلاسيكية \textbf{المتباعدة
بانتظام} (\(\lambda_n=n\)). أما تطبيق الكثافة فيحتاج المتباينة عند
\(\lambda_n=\log n\)، وهذه \textbf{غير منتظمة التباعد}
(\(\delta_n\asymp1/n\))، فتلزم الصيغة \textbf{الموزونة} لمونتغمري–فون
(القضية~\ref{prop:ch58-mv-weighted-cited})، ويلزم كذلك التقطيع إلى نقاط
متباعدة (القضية~\ref{prop:ch58-large-values-cited}). وكلتاهما مُستشهَد
بها هنا لا مبرهَنة. لذلك \textbf{لا يجوز} قراءة هذا الفصل بوصفه طريقًا
مكتملًا داخليًّا إلى حدود الكثافة.
\end{remark}

\subsection{المتطلبات السابقة}

لا يحتاج هذا الفصل مبرهنةً واحدة من المجلد الأول. السلسلة مكتفية بذاتها:
متسلسلة فورييه أولية واحدة (مُستشهَد بها)، وتكامل بالتجزئة، وتعامد
الدوال الأسّية.

\subsection{نواتج التعلم}

بعد إتمام الفصل ينبغي أن يستطيع القارئ:
\begin{enumerate}
  \item حساب معامِلات فورييه للدالة \(g(x)=i\pi(1-2x)\) من صيغتها المغلقة.
  \item إثبات التمثيل التكاملي للصيغة الثنائية لهيلبرت.
  \item استنتاج المتباينة بثابتها الحادّ \(\pi\).
  \item تفسير لماذا تتلاشى الصيغة عند المعاملات الحقيقية، ولماذا يلزم
  ذلك أن تكون أمثلة الحدّية مركّبة بالضرورة.
  \item التمييز بين ما أُثبت هنا وما استُشهد به، وحدود ما يمكن بناؤه على
  كلٍّ منهما.
\end{enumerate}

\section{الاصطلاح}
\symbolindex{الاصطلاح}

في كل ما يلي \(N\in\mathbb{N}\)، و\(z_1,\dots,z_N\in\mathbb{C}\) متتالية
\textbf{منتهية}، و
\[
e(y):=e^{2\pi iy},
\qquad
f(x):=\sum_{n=1}^{N}z_ne(nx).
\]
المجموع \(f\) منتهٍ، فلا تنشأ أي مسألة تقارب فيه. والمقام المعتمد في هذا
الفصل هو \(n-m\)؛ وصيغة \(m-n\) سالبها، فتتكافأ معها تحت القيمة المطلقة.

\section{النواة \texorpdfstring{\(g\)}{g} وحدّها الأعلى}
\symbolindex{النواة g وحدّها الأعلى}

نبدأ بحقيقة فورييه الكلاسيكية الوحيدة التي يستوردها الفصل.

\theoremindex{متسلسلة فورييه للدالة السنّية}
\begin{lemma}[متسلسلة الدالة السنّية]
\label{lem:ch58-sawtooth}
\resultid{ANT-LEM-58-01}
\citedresult{\cite{Ahlfors1979}}
لكل \(0<x<1\):
\[
\sum_{k\ge1}\frac{\sin(2\pi kx)}{k}=\pi\Bigl(\frac12-x\Bigr).
\]
\end{lemma}

هذه متسلسلة متقاربة تقاربًا شرطيًّا لا مطلقًا، وهي نقطة سنعود إليها
بعنايةٍ عند بناء التمثيل التكاملي.

\theoremindex{النواة g بصيغتها المغلقة}
\begin{lemma}[النواة \(g\) بصيغتها المغلقة]
\label{lem:ch58-kernel}
\resultid{ANT-LEM-58-02}
\provedhere
لكل \(0<x<1\):
\[
g(x):=\sum_{k\ne0}\frac{e(kx)}{k}=i\pi(1-2x),
\]
ومن ثَمّ \(\lvert g(x)\rvert<\pi\) لكل \(x\in(0,1)\).
\end{lemma}

\begin{proof}
نجمع الحدّين \(k\) و\(-k\) لكل \(k\ge1\):
\[
\frac{e(kx)}{k}+\frac{e(-kx)}{-k}
=\frac{e(kx)-e(-kx)}{k}
=\frac{2i\sin(2\pi kx)}{k}.
\]
إذن، بالجمع على \(k\ge1\) واستعمال اللمّة~\ref{lem:ch58-sawtooth}،
\[
g(x)=2i\sum_{k\ge1}\frac{\sin(2\pi kx)}{k}
=2i\cdot\pi\Bigl(\frac12-x\Bigr)
=i\pi(1-2x).
\]
أما الحدّ الأعلى: لكل \(x\in(0,1)\) لدينا \(-1<1-2x<1\)، أي
\(\lvert1-2x\rvert<1\)، فينتج \(\lvert g(x)\rvert=\pi\lvert1-2x\rvert<\pi\).
لاحظ أن المتراجحة \textbf{صارمة} داخل الفترة المفتوحة، ولا تبلغ \(\pi\)
إلا في النهايتين \(x\to0^{+}\) و\(x\to1^{-}\)، وهما مجموعة قياسها صفر.
\end{proof}

\section{معامِلات فورييه للنواة}
\symbolindex{معامِلات فورييه للنواة}

الخطوة التالية هي جوهر الفصل تقنيًّا. نحسب معامِلات فورييه لِـ\(g\)
\textbf{من صيغتها المغلقة} لا من متسلسلتها، فنتجنّب تمامًا الحاجة إلى
تبديل تكامل بمتسلسلة متقاربة تقاربًا شرطيًّا.

\theoremindex{تكاملان أوليان}
\begin{lemma}[تكاملان أوليان]
\label{lem:ch58-elementary-integrals}
\resultid{ANT-LEM-58-03}
\provedhere
لكل \(j\in\mathbb{Z}\setminus\{0\}\):
\[
\int_0^1e(jx)\,dx=0,
\qquad
\int_0^1x\,e(jx)\,dx=\frac{1}{2\pi ij}.
\]
\end{lemma}

\begin{proof}
الأول مباشر: \(\int_0^1e^{2\pi ijx}dx=\bigl[e^{2\pi ijx}/(2\pi ij)\bigr]_0^1=0\)
لأن \(e^{2\pi ij}=1\). والثاني بالتكامل بالتجزئة:
\[
\int_0^1x\,e^{2\pi ijx}\,dx
=\Bigl[\frac{x\,e^{2\pi ijx}}{2\pi ij}\Bigr]_0^1
-\frac{1}{2\pi ij}\int_0^1e^{2\pi ijx}\,dx
=\frac{1}{2\pi ij}-0
=\frac{1}{2\pi ij},
\]
حيث تلاشى التكامل الأخير بالنتيجة الأولى.
\end{proof}

\theoremindex{معامِلات فورييه للنواة g}
\begin{lemma}[معامِلات فورييه لِـ\(g\)]
\label{lem:ch58-kernel-coefficients}
\resultid{ANT-LEM-58-04}
\provedhere
لكل \(j\in\mathbb{Z}\):
\[
\int_0^1e(jx)\,g(x)\,dx=
\begin{cases}
-\dfrac1j, & j\ne0,\\[2mm]
0, & j=0.
\end{cases}
\]
\end{lemma}

\begin{proof}
نستعمل الصيغة المغلقة \(g(x)=i\pi(1-2x)\) من
اللمّة~\ref{lem:ch58-kernel}. عند \(j\ne0\)، وباللمّة~\ref{lem:ch58-elementary-integrals}:
\[
\int_0^1e(jx)\,i\pi(1-2x)\,dx
=i\pi\Bigl(\int_0^1e(jx)\,dx-2\int_0^1x\,e(jx)\,dx\Bigr)
=i\pi\Bigl(0-\frac{2}{2\pi ij}\Bigr)
=-\frac{i\pi}{\pi ij}
=-\frac1j.
\]
وعند \(j=0\):
\[
\int_0^1i\pi(1-2x)\,dx
=i\pi\Bigl[x-x^2\Bigr]_0^1
=i\pi(1-1)=0.
\]
\end{proof}

\begin{remark}[لماذا هذا المسار وليس تبديل المتسلسلة والتكامل]
\label{rem:ch58-why-closed-form}
الطريق الأقصر ظاهريًّا هو إدخال التكامل داخل متسلسلة \(g\) حدًّا حدًّا.
لكن تلك المتسلسلة متقاربة تقاربًا \textbf{شرطيًّا} لا مطلقًا
(اللمّة~\ref{lem:ch58-sawtooth})، فتبديل التكامل بها يحتاج مبرِّرًا
إضافيًّا (تقارب مهيمن أو حدٌّ منتظم على المجاميع الجزئية) لا نملكه هنا
مجانًا. أما الحساب أعلاه فيتكامل ضدّ \textbf{صيغة مغلقة} لا ضدّ متسلسلة،
فيسقط الإشكال من أصله. وستكون العملية الوحيدة المتبقية في
القضية~\ref{prop:ch58-integral-representation} تبديل تكامل بمجموع
\textbf{منتهٍ}، وهو تبديل مشروع دائمًا بلا شروط.
\end{remark}

\section{التمثيل التكاملي}
\symbolindex{التمثيل التكاملي}

\theoremindex{التمثيل التكاملي للصيغة الثنائية}
\begin{proposition}[التمثيل التكاملي]
\label{prop:ch58-integral-representation}
\resultid{ANT-PROP-58-01}
\provedhere
لكل \(z_1,\dots,z_N\in\mathbb{C}\):
\[
\int_0^1\lvert f(x)\rvert^2g(x)\,dx
=\sum_{\substack{m,n=1\\m\ne n}}^{N}\frac{z_m\overline{z_n}}{n-m}.
\]
\end{proposition}

\begin{proof}
بما أن \(f\) مجموع منتهٍ:
\[
\lvert f(x)\rvert^2
=f(x)\overline{f(x)}
=\sum_{m=1}^{N}\sum_{n=1}^{N}z_m\overline{z_n}\,e\bigl((m-n)x\bigr).
\]
هذا مجموع منتهٍ من دوال متصلة، فتبديله مع التكامل مشروع:
\[
\int_0^1\lvert f(x)\rvert^2g(x)\,dx
=\sum_{m=1}^{N}\sum_{n=1}^{N}z_m\overline{z_n}\int_0^1e\bigl((m-n)x\bigr)g(x)\,dx.
\]
نطبّق اللمّة~\ref{lem:ch58-kernel-coefficients} بوضع \(j=m-n\). الحدود
القطرية (\(m=n\)، أي \(j=0\)) تلاشت. وأما \(m\ne n\) فتعطي
\(-1/(m-n)=1/(n-m)\)، فينتج
\[
\int_0^1\lvert f(x)\rvert^2g(x)\,dx
=\sum_{m\ne n}z_m\overline{z_n}\cdot\frac{1}{n-m},
\]
وهو المطلوب.
\end{proof}

\section{متباينة هيلبرت بثابت شور}
\symbolindex{متباينة هيلبرت بثابت شور}

\theoremindex{متباينة هيلبرت بثابت شور الحادّ}
\begin{theorem}[متباينة هيلبرت، بثابت شور]
\label{thm:ch58-hilbert-inequality}
\resultid{ANT-THM-58-01}
\provedhere
لكل \(N\in\mathbb{N}\) وكل \(z_1,\dots,z_N\in\mathbb{C}\):
\[
\Biggl\lvert
\sum_{\substack{m,n=1\\m\ne n}}^{N}\frac{z_m\overline{z_n}}{n-m}
\Biggr\rvert
\le
\pi\sum_{n=1}^{N}\lvert z_n\rvert^2.
\]
\end{theorem}

\begin{proof}
من القضية~\ref{prop:ch58-integral-representation}، ثم بتقدير التكامل
بحدّ النواة الأعلى من اللمّة~\ref{lem:ch58-kernel}:
\[
\Biggl\lvert\sum_{m\ne n}\frac{z_m\overline{z_n}}{n-m}\Biggr\rvert
=\Biggl\lvert\int_0^1\lvert f(x)\rvert^2g(x)\,dx\Biggr\rvert
\le\sup_{x\in(0,1)}\lvert g(x)\rvert\cdot\int_0^1\lvert f(x)\rvert^2\,dx
\le\pi\int_0^1\lvert f(x)\rvert^2\,dx.
\]
ويبقى حساب \(\int_0^1\lvert f\rvert^2\). بفكّ المربع وتعامد الدوال
الأسّية (وهو نفس التعامد المستعمَل في
اللمّة~\ref{lem:ch58-elementary-integrals}):
\[
\int_0^1\lvert f(x)\rvert^2\,dx
=\sum_{m=1}^{N}\sum_{n=1}^{N}z_m\overline{z_n}\int_0^1e\bigl((m-n)x\bigr)dx
=\sum_{n=1}^{N}\lvert z_n\rvert^2,
\]
إذ لا يبقى إلا الحدّ القطري \(m=n\). بالتعويض ينتج المطلوب.
\end{proof}

\begin{remark}[لا حاجة إلى مبرهنة بارسيفال]
\label{rem:ch58-no-parseval}
يُغري الشكل الأخير باستدعاء مساواة بارسيفال، لكن ذلك تحميل زائد: المجموع
هنا \textbf{منتهٍ}، والمطلوب هو تعامد \(\{e(nx)\}\) على \([0,1]\) فحسب،
وهو ما استُعمل أصلًا في بناء التمثيل التكاملي. لا مبرهنة إضافية تدخل هنا.
\end{remark}

\section{ملاحظتان بنيويتان}
\symbolindex{ملاحظتان بنيويتان}

\theoremindex{تلاشي الصيغة عند المعاملات الحقيقية}
\begin{proposition}[تلاشي الصيغة عند المعاملات الحقيقية]
\label{prop:ch58-real-vanishing}
\resultid{ANT-PROP-58-02}
\provedhere
إذا كان \(z_n\in\mathbb{R}\) لكل \(n\)، فإن
\[
\sum_{\substack{m,n=1\\m\ne n}}^{N}\frac{z_mz_n}{n-m}=0.
\]
\end{proposition}

\begin{proof}
النواة \(K_{mn}:=1/(n-m)\) متخالفة التناظر: \(K_{nm}=-K_{mn}\). أما
الأوزان \(z_mz_n\) فمتناظرة عند المعاملات الحقيقية. إذن يتقابل كل حدّ
\((m,n)\) مع الحدّ \((n,m)\) فيُلغيه، ويتلاشى المجموع كلّه.
\end{proof}

\begin{remark}[أثر هذا التلاشي على البحث عن أمثلة حدّية]
\label{rem:ch58-real-vanishing-consequence}
هذه الملاحظة عملية لا تجميلية: تعني أن أي محاولة لاختبار حدّة الثابت
\(\pi\) بمتتاليات \textbf{حقيقية} ستقيس صفرًا لا أكثر، فتظهر نسبًا تقترب
من الصفر ويُستنتج خطأً أن الثابت بعيد عن الحدّية. المتتاليات الحدّية
لهذه الصيغة \textbf{مركّبة بالضرورة}.
\end{remark}

\theoremindex{صرامة المتباينة عند كل مقطع منته}
\begin{proposition}[صرامة المتباينة]
\label{prop:ch58-strictness}
\resultid{ANT-PROP-58-03}
\provedhere
لكل \(N\) ولكل \(z\ne0\)، المتراجحة في
المبرهنة~\ref{thm:ch58-hilbert-inequality} \textbf{صارمة}:
\[
\Biggl\lvert\sum_{m\ne n}\frac{z_m\overline{z_n}}{n-m}\Biggr\rvert
<\pi\sum_{n=1}^{N}\lvert z_n\rvert^2.
\]
\end{proposition}

\begin{proof}
من اللمّة~\ref{lem:ch58-kernel} لدينا \(\lvert g(x)\rvert<\pi\) صارمًا
لكل \(x\in(0,1)\). والدالة \(x\mapsto\lvert f(x)\rvert^2\) متصلة وغير
سالبة، ولا تتلاشى على \((0,1)\) كلّها (وإلا لكان \(f\equiv0\) ومن ثَمّ
\(z=0\) بتعامد الدوال الأسّية). إذن
\[
\Biggl\lvert\int_0^1\lvert f\rvert^2g\Biggr\rvert
\le\int_0^1\lvert f\rvert^2\lvert g\rvert
<\pi\int_0^1\lvert f\rvert^2,
\]
حيث جاءت الصرامة من أن \(\lvert g\rvert\) أصغر من \(\pi\) صرامةً على
مجموعة موجبة القياس يكون فيها \(\lvert f\rvert^2>0\).
\end{proof}

\begin{remark}[توفيق ظاهري مهم]
\label{rem:ch58-strictness-vs-sharpness}
لا يتعارض هذا مع حدّة الثابت \(\pi\)
(القضية~\ref{prop:ch58-sharpness-cited}). الحدّة تعني أن \(\pi\) هو
\textbf{أصغر} ثابت صالح لكل \(N\) معًا، أي أن نسبة الطرفين تقترب من
\(1\) عند \(N\to\infty\)؛ ولا تعني أنها تبلغ \(1\) عند أي \(N\) منتهٍ.
والصرامة أعلاه تقول بالضبط إنها لا تبلغه أبدًا عند \(N\) منتهٍ.
\end{remark}

\section{ما لا يُبرهَن هنا: النتائج المُستشهَد بها}
\symbolindex{النتائج المُستشهَد بها}

\theoremindex{حدّة ثابت شور}
\begin{proposition}[حدّة الثابت \(\pi\)]
\label{prop:ch58-sharpness-cited}
\resultid{ANT-PROP-58-04}
\citedresult{\cite{Schur1911,HardyLittlewoodPolya1952}}
الثابت \(\pi\) في المبرهنة~\ref{thm:ch58-hilbert-inequality} أمثل: لا
يمكن استبداله بأي ثابت أصغر يصلح لكل \(N\) ولكل المعاملات.
\end{proposition}

\begin{remark}[حالة هذه القضية بدقة]
\label{rem:ch58-sharpness-status}
\noindent\textenglish{CITED / NOT PROVED HERE}. لا يقدّم هذا الفصل
برهانًا للحدّة، ولا يقدّم بديلًا عنه. العرض المعياري في
\cite[المبرهنة~294]{HardyLittlewoodPolya1952}. وما يمكن قوله داخليًّا هو
أن حسابًا عدديًّا للمعيار المؤثِّري يعطي نسبًا متزايدة رتيبًا نحو
\(1\) دون بلوغه، وهو \textbf{شاهد متوافق} مع الحدّة لا برهان لها.
\end{remark}

\theoremindex{تعميم مونتغمري–فون الموزون}
\begin{proposition}[تعميم مونتغمري–فون الموزون]
\label{prop:ch58-mv-weighted-cited}
\resultid{ANT-PROP-58-05}
\citedresult{\cite{MontgomeryVaughan1974}}
لتكن \(\lambda_1<\dots<\lambda_N\) أعدادًا حقيقية، ولنضع
\(\delta_n:=\min\{\lambda_n-\lambda_{n-1},\ \lambda_{n+1}-\lambda_n\}\).
عندئذٍ يوجد ثابت مطلق \(C_1\) بحيث
\[
\Biggl\lvert\sum_{m\ne n}\frac{z_m\overline{z_n}}{\lambda_m-\lambda_n}\Biggr\rvert
\le C_1\sum_{n=1}^{N}\frac{\lvert z_n\rvert^2}{\delta_n}.
\]
\end{proposition}

\begin{remark}[الفارق الجوهري عن المبرهنة المُثبتة أعلاه]
\label{rem:ch58-mv-vs-schur}
المبرهنة~\ref{thm:ch58-hilbert-inequality} مقصورة على التباعد المنتظم
(\(\lambda_n=n\)، فتباعدها \(1\)). والقضية أعلاه تسمح بتباعد
\textbf{غير منتظم} عبر أوزان \(\delta_n\) لكل حدّ على حدة — وهذا بالضبط
ما يلزم في التطبيقات، حيث تكون \(\lambda_n=\log n\) فيصير
\(\delta_n\asymp1/n\). الانتقال بين الحالتين ليس تفصيلًا تقنيًّا بل هو
موضع العمق، ولذلك لا يُبرهَن هنا.

الثابت الذي أثبته مونتغمري وفون صريح لكن \textbf{غير أمثل}؛ وتحديد
الثابت الأمثل مسألة \textbf{مفتوحة} إلى اليوم. ويجب عدم الخلط بين
الأمرين: المتراجحة نفسها \textbf{مبرهَنة ومغلقة}، والمفتوح هو تحسين
ثابتها فقط. وهذا التمييز كافٍ تمامًا لأي استعمال بصيغة \(\ll\).
\end{remark}

\theoremindex{مبرهنة القيم الكبيرة لكثيرات حدود ديريشليه}
\begin{proposition}[القيم الكبيرة لكثيرات حدود ديريشليه]
\label{prop:ch58-large-values-cited}
\resultid{ANT-PROP-58-06}
\citedresult{\cite{Montgomery1969}}
ليكن \(F(s)=\sum_{n\le N}a_nn^{-s}\)، ولتكن \(t_1,\dots,t_R\) نقاطًا في
\([-T,T]\) متباعدة بمسافة \(\ge1\). إذا كان \(\lvert F(it_r)\rvert\ge V\)
لكل \(r\)، فإن
\[
R\ll\frac{(N+T)\log(2N)}{V^2}\sum_{n\le N}\lvert a_n\rvert^2.
\]
\end{proposition}

\section{ما الذي لم يثبته هذا الفصل؟}
\symbolindex{ما الذي لم يثبته هذا الفصل؟}

\begin{itemize}
  \item لم يثبت حدّة الثابت \(\pi\)
  (القضية~\ref{prop:ch58-sharpness-cited}) — مُستشهَد بها.
  \item لم يثبت تعميم مونتغمري–فون الموزون
  (القضية~\ref{prop:ch58-mv-weighted-cited}) — مُستشهَد به، وهو المدخل
  الفعلي لأي تطبيق على الكثافة.
  \item لم يثبت مبرهنة القيم الكبيرة
  (القضية~\ref{prop:ch58-large-values-cited}) — مُستشهَد بها.
  \item لم يشتق أي حدّ كثافة صفرية، ولا أي أسّ عددي. راجع التحذير في
  الملاحظة~\ref{rem:ch58-scope-warning}.
\end{itemize}

\section{تمارين}
\symbolindex{تمارين}

\begin{enumerate}
  \item تحقّق مباشرةً عند \(N=2\) و\(z=(1,i)\) أن طرفَي التمثيل التكاملي
  في القضية~\ref{prop:ch58-integral-representation} متساويان، وأن
  قيمتهما \(-2i\).
  \item بيّن أن الصيغة الثنائية لهيلبرت \(\sum_{m\ne n}z_m\overline{z_n}/(n-m)\)
  تخيّلية محضة لكل \(z\). (تلميح: النواة متخالفة التناظر، فما هو مرافق
  الصيغة؟)
  \item اشرح لماذا لا ينفع اختيار \(z_n=n^{-1/2}\) لاختبار حدّة \(\pi\)
  عدديًّا، بالرجوع إلى القضية~\ref{prop:ch58-real-vanishing}.
  \item أين بالضبط يفشل برهان المبرهنة~\ref{thm:ch58-hilbert-inequality}
  لو استُبدل \(\lambda_n=n\) بـ\(\lambda_n=\log n\)؟ حدِّد الخطوة
  بعينها، ثم قارن بالملاحظة~\ref{rem:ch58-mv-vs-schur}.
\end{enumerate}

\section{خلاصة الفصل}
\symbolindex{خلاصة الفصل}

أثبت هذا الفصل متباينة هيلبرت بثابت شور الحادّ \(\pi\) عبر سلسلة قصيرة
مغلقة: متسلسلة سنّية مُستشهَد بها، ثم صيغة مغلقة للنواة \(g\) وحدّ أعلى
صارم لها، ثم معامِلات فورييه للنواة محسوبة من الصيغة المغلقة (فلا حاجة
إلى تبديل تكامل بمتسلسلة متقاربة شرطيًّا)، ثم تمثيل تكاملي للصيغة
الثنائية لا يتضمن إلا تبديل مجموع منتهٍ، ثم المتباينة. وأُضيفت ملاحظتان
بنيويتان: تلاشي الصيغة عند المعاملات الحقيقية، وصرامة المتراجحة عند كل
\(N\) منتهٍ.

أما الأدوات التي تُستعمَل فعلًا في مسائل الكثافة الصفرية — الصيغة
الموزونة لتباعد غير منتظم، ومبرهنة القيم الكبيرة — فقد عُرضت
\textbf{مُستشهَدًا بها} صراحةً. الفارق بين التباعد المنتظم وغير المنتظم
ليس تفصيلًا يُطوى، بل هو الحدّ الفاصل بين ما يمكن بناؤه هنا وما يحتاج
آلةً أثقل. وقد نُصَّ على ذلك في صدر الفصل لئلا يُقرأ ما بعده على أنه
طريق مكتمل.
